% ======================================================================
% Ontology of Continua — Core 1.3
% Cross-K module: crossk_k0_k1.tex
% Cross-level relations between K0 and K1
% ======================================================================

\subsubsection{Межуровневая структура K0–K1}
\label{sec:crossk-k0-k1}

Переход от $K_0$ к $K_1$ представляет собой первую и наиболее фундаментальную
межуровневую структуру во всей иерархии континуумов. Он описывает
появление различимых состояний, возникновение явной оси и формирование
нетривиального пространства состояний $\Omega(K_1)$ из минимальной
предонтологической основы $K_0$.

В этом разделе суммируются структурные зависимости, наследуемые компоненты,
операторы перехода и пороговые условия возникновения $K_1$.

% ----------------------------------------------------------------------
\subsubsection{1. Задействованные уровни и их роли}

\paragraph{K0.}
$K_0$ определяется \emph{аксиомой различия и связности}.
Он содержит:
\begin{itemize}
  \item минимальное множество потенциальных состояний без явных осей;
  \item структурную меру различия $\Delta(s_1,s_2)$;
  \item порог существования $\Theta_0$, обеспечивающий нетривиальность:
        $0 < \Delta < \varepsilon$ при некотором $\varepsilon>0$;
  \item без времени, без потоков $J$, без циклов $C$, без динамики.
\end{itemize}

\paragraph{K1.}
$K_1$ — первый континуум с:
\begin{itemize}
  \item явной осью $A_1$, захватывающей устойчивый класс различий;
  \item определенным пространством состояний $\Omega(K_1)$ с границей
        $\partial\Omega(K_1)$;
  \item примитивной энергопотенциальной структурой $P_1$;
  \item минимальными потоками $J_1$ и тривиальным циклом
        $C_{\mathrm{triv}}$;
  \item явными порогами $\Theta_1$ (существования, градиента, устойчивости).
\end{itemize}

Роль перехода K0→K1 — это \textbf{генерация структуры}: формирование
первого представимого измерения в континууме.

% ----------------------------------------------------------------------
\subsubsection{2. Совместные и наследуемые оси и пороги}

Хотя у $K_0$ нет осей, несколько его структурных свойств наследуются
в $K_1$ в преобразованном виде:

\paragraph{Различие \texorpdfstring{$\Delta$}{\Delta}.}
Примитивная структура различий $K_0$ становится метрическим компонентом оси
$A_1$ в $K_1$:
\[
A_1 \sim f(\Delta),
\]
то есть существование различимости в $K_0$ и есть то, что позволяет
появиться явной координате в $K_1$.

\paragraph{Пороги.}
Порог существования $\Theta_0$ становится \emph{нижней границей} для $K_1$:
если $\Theta_0$ не выполняется (то есть $\Delta=0$), ось не может быть
сформирована и $\Omega(K_1)=\varnothing$.

Напротив, $\Theta_1$ вводит новые условия устойчивости для градиентов,
границ и допустимых потоков. Эти пороги отсутствуют в $K_0$ и
возникают только после появления первой оси.

% ----------------------------------------------------------------------
\subsubsection{3. Межуровневые потоки, циклы и напряжения}

\paragraph{Потоки.}
В $K_0$ нет внутренних потоков, поэтому оператор $J_1$ уровня $K_1$ не имеет
предшественника. Вместо этого он возникает из только что доступной
структуры:
\[
J_1 : A_1 \rightarrow \Omega(K_1)
\]
представляет движения вдоль или поперёк новой оси.

\paragraph{Циклы.}
$K_0$ не поддерживает циклы, поскольку в нём отсутствуют время и потоки.
Поэтому тривиальный цикл $K_1$:
\[
C_{\mathrm{triv}} : s \mapsto s,
\]
является первой возможной циклической структурой в иерархии.

\paragraph{Напряжение.}
Межуровневое структурное напряжение между $K_0$ и $K_1$ определяется как:
\[
T_{0,1} = g(\Delta - \Theta_0,\ P_1 - \Theta_1).
\]
Если $T_{0,1}$ превышает его размерностный порог $\Theta_{1,\mathrm{dim}}$,
возникшая ось становится неустойчивой или схлопывается.

% ----------------------------------------------------------------------
\subsubsection{4. Условия рождения и исчезновения между уровнями}

\paragraph{Рождение \texorpdfstring{$K_1$}{K_1}.}
$K_1$ возникает при:
\[
\Delta > \Theta_0
\quad \text{и} \quad
T_0 > \Theta_{0,\mathrm{dim}},
\]
что позволяет появиться новой оси $A_1$:
\[
A_1 \in M_1 \setminus A(K_0).
\]

Это действие оператора $\Psi_{0\rightarrow 1}$ формализовано в Core:
он расширяет минимальную структуру $K_0$ в представительный континуум
$K_1$.

\paragraph{Распространение гибели.}
Если $\Theta_0$ не выполняется, $K_1$ не может существовать:
\[
\Theta_0^{\mathrm{death}}
\Rightarrow \Omega(K_1)=\varnothing.
\]

Если $\Theta_1$ не выполняется, континуум $K_1$ схлопывается, но $K_0$
остаётся целым, поскольку у $K_0$ нет структурной зависимости от $K_1$.

% ----------------------------------------------------------------------
\subsubsection{5. Концептуальные примеры}

\paragraph{Возникновение первой оси из минимальных различий.}
Любая ситуация, где необработанное различие (например, «состояние A отличается
от состояния B») может быть стабилизировано и параметризовано, соответствует
реализации перехода K0→K1. Ось $A_1$ воплощает это стабилизированное
различие.

\paragraph{Формирование границ.}
С появлением оси границы обретают смысл:
\[
\partial\Omega(K_1) = \{ s \in \Omega(K_1) \mid \nabla A_1 \text{ достигает порога} \}.
\]

\paragraph{Первые потоки.}
При наличии оси и границ возникают градиентные подобные потоки:
возможность, невозможная в $K_0$, но естественное следствие стабилизированного
различия в $K_1$.

Эти примеры иллюстрируют глобальное значение перехода K0→K1:
\[
\textbf{Это рождение всей представимой структуры в иерархии континуумов.}
\]
